\chapter{国内外研究现状}\label{chap:related}

人工智能辅助多项式时间归约构造涉及归约的表达与验证、数学证明生成、NP 问题实例求解和智能体交互。现有研究已经分别处理了这些环节，但各自采用的研究对象与正确性标准并不相同。

\section{NP 困难性证明的生成与验证}

Creus 等使用 SAT 方法评价候选归约，实验中发现了纸笔构造的具体错误\citep{creus2014automatic}。这类评价依赖可编码的候选和实例，所得反例不能证明归约对所有输入成立。Zhang 等提出领域专用语言 Karp，用统一语义描述问题、实例映射和证书映射\citep{zhang2022karp}。该语言通过代价语义约束运行时间，并用性质测试检查了 25 个教材归约；受限语法排除了部分错误，也限制了可表达的归约类型。

Grange 等研究候选归约的自动验证问题，证明部分受限表示下的验证仍不可判定\citep{grange2024verification}。他们提出图形化的 cookbook reduction，并给出若干可判定子类。可判定性依赖表示形式的限制，不能直接推广到通用程序。

Karp 数据集收集了 90 个自然语言 NP 困难性归约证明，用于训练和评价语言模型\citep{dicicco2025karpdataset}。该数据集覆盖教材练习和论文中的归约，使源问题选择、映射构造和双向论证可在同一组材料上考察。自然语言证明中的边界条件和量词未必全部写明，评价仍较多依赖理论计算机科学专家。

Pan 等利用问题约束、知识库和多层检查构建可执行归约库，并通过归约关系图连接问题与已有求解器\citep{pan2026reductions}。结果表明代码智能体适于实现和维护成规模的已知归约，但其研究对象不包括复杂度未知问题中的源问题选择和新映射发现。

\section{人工智能辅助数学证明}

miniF2F 收集了 488 道形式化数学题\citep{zheng2022minif2f}。数据覆盖 Lean、Metamath、Isabelle 和 HOL Light 等系统。它为跨系统评价提供了统一材料，但形式化陈述已经给定。ProofNet 将本科数学教材中的自然语言命题、证明和 Lean 陈述配对，用于研究自动形式化与证明生成\citep{azerbayev2023proofnet}。其评价对象仍是预先选定的定理，不涉及待证结论的发现。


LeanDojo 从 Lean 项目中提取证明步骤与环境信息，并据此构建检索增强的定理证明器\citep{yang2023leandojo}。其结果说明证明库检索有助于前提选择，但检索范围受已有形式化库限制。Jiang 等将非形式证明转化为形式草图，再由自动证明器补全子目标\citep{jiang2023draft}。该方法证明非形式推理可引导形式证明搜索，但仍以给定定理和可调用的自动证明器为前提。

AlphaGeometry 用语言模型引导符号演绎程序，并通过合成数据学习欧氏几何证明\citep{trinh2024alphageometry}。其几何证明由符号程序检查，适用范围受几何语言和推理规则限制。AlphaProof 在 Lean 环境中通过强化学习搜索形式证明，完成了部分国际数学奥林匹克竞赛题\citep{hubert2026alphaproof}。逻辑内核保证了证明的正确性，但题目定义和待证结论已经固定。

\section{NP 问题实例求解与判定反馈}

NPHardEval 按复杂度类别组织 900 个算法问题，并通过动态更新题目减少测试数据泄漏\citep{fan2024nphardeval}。实验刻画了语言模型性能随问题难度和实例规模的变化，但评价单位仍是单个输入实例。

NPPC 提供 25 类 NP 完全问题的实例生成、求解和分析接口\citep{yang2025nppc}。持续生成的实例与答案检查程序支持规模可控的自动评价，但不产生适用于全部实例的映射函数。NP-Engine 将实例生成器、答案检查程序和启发式求解器结合，用于评价语言模型求解 NP 困难优化问题的表现\citep{li2025npengine}。合成实例的输出可自动检查，其结论仍限于给定实例上的可行解和目标函数值。

\section{代码智能体与运行评测框架}

ReAct 使语言模型交替生成文本推理和外部操作，并根据工具结果更新后续步骤\citep{yao2023react}。该方法表明语言模型可利用多轮环境反馈，但反馈本身不保证推理结论正确。

\begin{figure}[!htbp]
    \centering
    \includegraphics[width=0.82\textwidth]{react-agent-architecture-zh}
    \bicaption{ReAct 智能体的基本交互结构。}{Basic interaction structure of a ReAct agent.}
    \label{fig:react-agent-architecture}
\end{figure}

SWE-bench 以真实代码库中的软件问题评价模型生成补丁的能力\citep{jimenez2024swebench}。测试结果为补丁提供自动反馈，但任务通常已有测试或历史修复作为参照。SWE-agent 支持智能体浏览代码库、编辑文件和运行测试\citep{yang2024sweagent}。OpenHands 提供终端执行、网页浏览和受限运行环境\citep{wang2024openhands}。这些系统使智能体得以保留中间文件并根据执行结果多次修改程序，但测试通过不能替代程序对全部输入成立的证明。

智能体运行与评测框架把复杂任务所需的工具、状态和反馈组织为连续交互。证明探索中的反馈包括候选程序的运行信息、输入格式检查、实例判定结果和证明条件的核验状态。该框架通过提供这些信息，复现研究者在源问题选择、映射构造和证明整理之间往返的过程，使人工智能辅助发现新的 NP 完全问题成为可能。
